%% IJIET submission draft (IEEE two-column style, matching published IJIET layout).
%% Source of empirical numbers: four-partition learner-based summaries in analysis/four_partition/
%% and seed-42 gate tables in analysis/direction_c/.
%% Do not submit while the JEDM manuscript is under review (dual submission).
\documentclass[journal,10pt,twocolumn]{IEEEtran}

\usepackage{amsmath,amssymb}
\usepackage{graphicx}
\usepackage{booktabs}
\usepackage{multirow}
\usepackage{array}
\usepackage{url}
\usepackage[hidelinks]{hyperref}
\usepackage{cite}

\hyphenation{Know-ledge ASSISTments cal-i-bra-tion}

\begin{document}

\title{Sparse-Concept Calibration of Knowledge Tracing Models\\
for Mastery and Remediation Gates}

\author{Khanh-Trinh~Nguyen,
Tuan~Dao~Minh,
Duong~Nguyen~Tien,
Chi~Thanh~Nguyen,
and~Van-Hau~Nguyen%
\thanks{K.-T. Nguyen, T. Dao Minh, D. Nguyen Tien, and V.-H. Nguyen are with Hung Yen University of Technology and Education, Hung Yen, Vietnam (e-mail: trinhnk@utehy.edu.vn; tuanymc@utehy.edu.vn; duongnt@utehy.edu.vn; haunv@utehy.edu.vn).}%
\thanks{C. T. Nguyen is with the Academy of Military Science and Technology, Ha Noi, Vietnam (e-mail: thanhnc@ioit.ai.vn).}%
\thanks{Corresponding author: Van-Hau Nguyen (e-mail: haunv@utehy.edu.vn).}}

\markboth{International Journal of Information and Education Technology}%
{Nguyen \MakeLowercase{\textit{et al.}}: Sparse-Concept Calibration of Knowledge Tracing}

\maketitle

\begin{abstract}
Knowledge Tracing (KT) systems that drive practice, remediation, or advancement typically consume a predicted probability, not an area-under-the-curve (AUC) score. Aggregate ranking metrics can look acceptable while those probabilities are poorly calibrated on rarely practiced skills. This paper reports a diagnostic evaluation---not a new KT architecture---of how calibration behaves across train-only knowledge-component (KC) frequency strata on three public logs (ASSISTments 2012, Junyi Academy, and XES3G5M) with IRT, DKT, and SimpleKT. Sparse training evidence does not universally reduce discrimination. It can, however, expose a dataset-dependent calibration vulnerability. On ASSISTments 2012, SimpleKT expected calibration error (ECE) rises from $0.114$ on dense KCs to $0.228$ on sparse KCs (Limited occupancy, $N=415$); Junyi has no learner-based sparse stratum, and XES3G5M SimpleKT ECE is essentially flat. Treating the same probabilities as a locked global mastery gate at $\tau=0.7$ raises SimpleKT false mastery among advances from $0.196$ (dense) to $0.268$ (sparse) on one ASSISTments fold; the sparse--dense gap stays positive on all five training seeds (mean $0.047$; four unique student partitions). A train-only graph KT instantiation shrinks that seed-42 gap. The gate is a simulation, not a classroom trial. The practical implication for educational-technology designers is to validate thresholded decisions per frequency stratum rather than on population AUC or ECE alone.
\end{abstract}

\begin{IEEEkeywords}
Knowledge tracing, learning analytics, calibration, sparse concepts, mastery threshold, educational decision support.
\end{IEEEkeywords}

\section{Introduction}

Adaptive practice engines and intelligent tutoring systems increasingly use Knowledge Tracing (KT) to estimate whether a learner has mastered a skill and to decide whether to skip, remediate, or advance~\cite{corbett1994knowledge,piech2015deep,abdelrahman2023knowledge}. In deployment, that decision is rarely the AUC of a leaderboard. It is a probability $p$ compared with a threshold $\tau$: if $p$ is high enough, the system treats the skill as mastered and withholds extra help.

KT papers typically report population AUC and accuracy~\cite{liu2023simplekt,liu2022pykt}. Those metrics are useful for ranking models, but they can hide reliability problems on low-frequency knowledge components (KCs)---concepts with little training-fold evidence---because most test events sit on dense, frequently practiced skills. A model can rank well overall and still be overconfident on the sparse tail. When $p$ is used as a mastery or remediation gate, that overconfidence is not a cosmetic reporting issue: it is a false-mastery error.

This paper treats that mismatch as an evaluation problem for educational technology, not as a reason to propose another KT backbone. We ask three questions that a system designer can act on:

\begin{enumerate}
\item Does lower KC training frequency systematically worsen ranking (AUC) across datasets?
\item Does it worsen probability reliability (calibration), and is that pattern dataset-dependent?
\item If a global threshold is applied to $p$, does the sparse stratum produce a dirtier advance decision than the dense stratum?
\end{enumerate}

The empirical story is bounded. Sparse training evidence does \emph{not} universally degrade discrimination. On XES3G5M, sparse AUC is higher than dense AUC for DKT and SimpleKT. Calibration, however, can be more sensitive than aggregate ranking on some logs. On ASSISTments 2012, SimpleKT ECE increases from dense to sparse concepts, and a locked gate at $\tau=0.7$ produces a higher false-mastery rate among sparse advances. The same ECE gradient is absent on Junyi (empty sparse bucket) and essentially absent for SimpleKT on XES3G5M. We therefore do not claim that ``sparse KCs always fail.'' We claim that occupancy-aware, stratum-wise calibration---and, when probabilities are consumed as gates, a simulated decision error---should be checked before a KT model is wired into remediation logic.

Contributions, stated conservatively:

\begin{itemize}
\item A train-only KC-frequency protocol with an explicit strict cold-start group, so sparse diagnostics cannot leak test-fold counts into the definition of ``sparse.''
\item Per-stratum calibration (ECE, Brier decomposition) on three public datasets, with sample-size flags (Reliable / Limited / Insufficient).
\item A locked-threshold simulation of false mastery and miss rates, plus a five-seed check of the sparse--dense false-mastery gap on ASSISTments 2012.
\end{itemize}

We do not claim a new audit theory, a new calibrator, or a classroom RCT. Graph KT (GKT) and a CL4KT-style adapter appear only as a single-fold diagnostic instantiate on ASSISTments, not as a model-contribution bake-off.

\section{Related Work}

Classical KT uses Bayesian knowledge tracing or item-response models~\cite{corbett1994knowledge,rasch1960probabilistic}. Deep sequential models (DKT, attentive and transformer KT) improved aggregate prediction~\cite{piech2015deep,ghosh2020context,liu2023simplekt}. Graph and contrastive variants often motivate architecture by sparse or related skills~\cite{nakagawa2019graph,lee2022cl4kt}. Those papers still report mainly population metrics.

Educational measurement has long warned that ranking is not reliability~\cite{pelanek2015metrics}. Calibration error and reliability diagrams quantify whether predicted probabilities match empirical correctness~\cite{guo2017calibration,naeini2015obtaining,brier1950verification,degroot1983comparison}. Recent KT work has begun to study calibration and sparse-concept evaluation~\cite{yan2026recovering,huang2023sparsekt,bhattacharjee2025coldstart}, but deployed tutoring systems still typically consume a single global $p$. This study sits in that gap: we keep standard models, stratify by training-only frequency, and ask whether a population gate would treat sparse skills more leniently---or more dangerously---than dense ones.

\section{Method}

\subsection{Datasets, splits, and models}

We use three de-identified public logs after a common schema: ASSISTments 2012~\cite{assistments2012data,liu2022pykt}, Junyi Academy~\cite{junyi2019kaggle}, and XES3G5M~\cite{liu2023xes3g5m}. Table~\ref{tab:data} summarizes processed cohort size. Learner-based splits (unseen students) are the primary setting. Temporal splits are a complementary stress test and are not the source of the gate numbers below.

\begin{table}[!t]
\caption{Processed cohort statistics (learner-based split).}
\label{tab:data}
\centering
\footnotesize
\begin{tabular}{lrrrr}
\toprule
Dataset & Learners & KCs & Interactions & Test events \\
\midrule
ASSISTments 2012 & 27{,}806 & 265 & 2.66M & 534{,}150 \\
Junyi Academy & 71{,}014 & 1{,}326 & 16.2M & 3{,}269{,}022 \\
XES3G5M & 18{,}066 & 866 & 7.95M & 1{,}589{,}145 \\
\bottomrule
\end{tabular}
\end{table}

Baselines: one-parameter IRT~\cite{rasch1960probabilistic} as a classical reference, DKT~\cite{piech2015deep}, and SimpleKT~\cite{liu2023simplekt}. On ASSISTments fold~0 only we also score a train-only GKT graph~\cite{nakagawa2019graph} and a CL4KT protocol adapter~\cite{lee2022cl4kt} (contrastive views on \emph{training} sequences; not an official CL4KT checkpoint). IRT under learner-based splits has no ability parameter for unseen students, so its AUC is $0.50$ by construction; we report it as a base-rate reference, not as a ranking competitor.

Deep models are trained at seeds $42, 2024, 2025, 2026, 2027$. The processed partitions for seeds 2025 and 2026 are identical (\texttt{fold\_2}$=$\texttt{fold\_3}) on all three datasets. Tables that report mean$\pm$sd therefore summarize \emph{four unique student partitions}: the two initializations on the duplicated split are averaged first. The gate-robustness table still lists all five training seeds without dropping one after the fact.

\subsection{Train-only frequency strata}

For each fold, KC frequency $f_{\mathrm{train}}$ is counted on the training file only. Buckets: strict cold-start ($f=0$), very sparse ($0<f<20$), sparse ($20\le f<100$), medium ($100\le f<500$), dense ($f\ge 500$). Fig.~\ref{fig:strata} shows that dense KCs dominate event volume while a non-empty sparse-like tail exists on ASSISTments and XES3G5M. Occupancy flags follow test-event count $N$: Reliable ($N\ge 1000$), Limited ($100\le N<1000$), Insufficient ($N<100$). Success claims require Limited or Reliable occupancy. Very-sparse ASSISTments cells are Insufficient and descriptive only.

\begin{figure}[!t]
\centering
\includegraphics[width=\columnwidth]{../paper/figures/figure2_bucket_distribution.pdf}
\caption{Train-only KC-frequency strata. Dense concepts dominate interactions; sparse and cold-start KCs remain the diagnostic tail.}
\label{fig:strata}
\end{figure}

\subsection{Calibration}

Let $y\in\{0,1\}$ be correctness and $p\in[0,1]$ the predicted probability. With $M=15$ equal-width bins, ECE is
\begin{equation}
\mathrm{ECE}=\sum_{m=1}^{M}\frac{n_m}{N}\bigl|\mathrm{acc}_m-\mathrm{conf}_m\bigr|.
\end{equation}
We also report the Brier score and its reliability / resolution / uncertainty decomposition~\cite{degroot1983comparison}. Lower ECE is better calibration; it is not a substitute for AUC.

\subsection{Simulated mastery gate}

For a locked global threshold $\tau$, the system \emph{advances} (skips remediation) if $p\ge\tau$ and \emph{triggers remediation} otherwise. Sparse events never receive their own tuned $\tau$. The display threshold is $\tau=0.7$; the grid $\{0.5,0.6,0.7,0.8\}$ is recorded in the artifact and is not a search over sparse error.

Let $A=\{p\ge\tau\}$. False mastery (FM) is $P(y=0\mid A)$: among advances, how often was the answer wrong. Miss is $P(A\mid y=0)$: among wrong answers, how often the system still skipped help. If the model were calibrated on the advance set, FM would match $E[1-p\mid A]$. We report
\begin{equation}
\Delta\mathrm{FM}=\mathrm{FM}_{\mathrm{sparse}}-\mathrm{FM}_{\mathrm{dense}}.
\end{equation}
Positive $\Delta\mathrm{FM}$ means sparse advances are dirtier than dense advances. This is a simulated decision error, not an instructional RCT.

\section{Results}

\subsection{Aggregate ranking is not the sparse story}

Table~\ref{tab:overall} shows learner-based AUC. DKT is slightly above SimpleKT on all three datasets. IRT AUC is $0.50$ for the reason given above. These numbers would look like a routine KT bake-off. They do not tell a designer whether $p$ is trustworthy on the tail.

\begin{table}[!t]
\caption{Overall learner-based AUC (mean$\pm$sd over four unique partitions).}
\label{tab:overall}
\centering
\footnotesize
\begin{tabular}{llcc}
\toprule
Dataset & Model & AUC & ACC \\
\midrule
ASSISTments 2012 & IRT & $0.5000$ & $0.6973\pm0.0004$ \\
 & DKT & $0.6979\pm0.0014$ & $0.7182\pm0.0014$ \\
 & SimpleKT & $0.6837\pm0.0025$ & $0.6996\pm0.0032$ \\
\midrule
Junyi Academy & IRT & $0.5000$ & $0.7053\pm0.0018$ \\
 & DKT & $0.7320\pm0.0009$ & $0.7343\pm0.0015$ \\
 & SimpleKT & $0.7231\pm0.0030$ & $0.7274\pm0.0015$ \\
\midrule
XES3G5M & IRT & $0.5000$ & $0.7961\pm0.0031$ \\
 & DKT & $0.8171\pm0.0022$ & $0.8327\pm0.0032$ \\
 & SimpleKT & $0.7557\pm0.0013$ & $0.8067\pm0.0037$ \\
\bottomrule
\end{tabular}
\end{table}

On XES3G5M, sparse AUC is \emph{higher} than dense AUC (DKT $0.857$ vs.\ $0.817$; SimpleKT $0.847$ vs.\ $0.755$; Reliable occupancy). Lower frequency is therefore not a universal ranking failure. Junyi's learner-based sparse bucket is empty under the pre-registered cuts---a protocol outcome, not a missing cell to impute.

\subsection{Calibration can move when ranking does not}

Table~\ref{tab:ece} is the calibration punchline. On ASSISTments 2012, SimpleKT ECE increases from $0.1136\pm0.0066$ (dense, Reliable, $N=523{,}971$) to $0.1541$ (medium) to $0.2280\pm0.0197$ (sparse, Limited, $N=415$). DKT shows the same direction (dense ECE $0.060$ vs.\ sparse $0.233$). IRT dense ECE is tiny ($0.003$) but resolution is zero: a constant-like predictor can look ``calibrated'' while ranking nothing.

On Junyi, only dense and medium strata exist; SimpleKT ECE rises modestly from $0.079$ to $0.107$. On XES3G5M, SimpleKT ECE is essentially flat ($0.114\to 0.111\to 0.125$ dense/medium/sparse) despite Reliable sparse occupancy ($N=2{,}010$). ECE-flat is not a license to skip occupancy reporting, and---as the gate results show---it is not the same as a flat miss rate.

\begin{table}[!t]
\caption{SimpleKT event-level ECE by train-only frequency stratum (four unique partitions). Flags: R Reliable, L Limited. Junyi sparse is empty.}
\label{tab:ece}
\centering
\footnotesize
\begin{tabular}{llrcc}
\toprule
Dataset & Stratum & $N$ & Flag & ECE \\
\midrule
ASSISTments 2012 & dense & 523{,}971 & R & $0.1136\pm0.0066$ \\
 & medium & 5{,}963 & R & $0.1541\pm0.0051$ \\
 & sparse & 415 & L & $0.2280\pm0.0197$ \\
\midrule
Junyi Academy & dense & 3{,}232{,}614 & R & $0.0792\pm0.0051$ \\
 & medium & 3{,}836 & R & $0.1073\pm0.0156$ \\
\midrule
XES3G5M & dense & 1{,}268{,}696 & R & $0.1145\pm0.0011$ \\
 & medium & 12{,}980 & R & $0.1114\pm0.0076$ \\
 & sparse & 2{,}010 & R & $0.1248\pm0.0085$ \\
\bottomrule
\end{tabular}
\end{table}

\subsection{A global gate can turn the ASSISTments gradient into a decision error}

Table~\ref{tab:gate} applies $\tau=0.7$ on ASSISTments fold~0 (seed 42; sparse $N=444$, Limited). SimpleKT false mastery among advances is $0.196$ on dense KCs and $0.268$ on sparse KCs ($\Delta\mathrm{FM}=+0.072$). The calibrated expectation $E[\mathrm{FM}]$ on sparse advances is only $0.050$, so the excess error is large. DKT shows a similar FM gap. Train-only GKT shrinks $\Delta\mathrm{FM}$ to $+0.015$ (FM $0.205\to 0.220$). The CL4KT adapter is intermediate ($0.185\to 0.240$). We do not read this as ``use GKT in production''; we read it as: the ASSISTments calibration gradient is not a property of the dataset alone.

\begin{table}[!t]
\caption{Simulated gate at $\tau=0.7$, ASSISTments 2012 fold 0. Advance if $p\ge\tau$. FM$=P(y=0\mid p\ge\tau)$; $E[\mathrm{FM}]=E[1-p\mid A]$. Not a classroom trial.}
\label{tab:gate}
\centering
\setlength{\tabcolsep}{3.5pt}
\begin{tabular}{llrrrr}
\toprule
Model & Stratum & $N$ & FM & $E[\mathrm{FM}]$ & Miss \\
\midrule
SimpleKT & dense & 528{,}018 & 0.196 & 0.113 & 0.352 \\
SimpleKT & sparse & 444 & 0.268 & 0.050 & 0.320 \\
DKT & dense & 528{,}018 & 0.200 & 0.147 & 0.398 \\
DKT & sparse & 444 & 0.296 & 0.057 & 0.365 \\
GKT (train-only) & dense & 528{,}018 & 0.205 & 0.163 & 0.455 \\
GKT (train-only) & sparse & 444 & 0.220 & 0.157 & 0.234 \\
CL4KT (adapter) & dense & 528{,}018 & 0.185 & 0.175 & 0.359 \\
CL4KT (adapter) & sparse & 444 & 0.240 & 0.116 & 0.244 \\
\bottomrule
\end{tabular}
\end{table}

Table~\ref{tab:five} checks whether SimpleKT $\Delta\mathrm{FM}$ is a one-fold accident. It is positive on all five training seeds (mean $0.047$, sd $0.033$). A KC-clustered bootstrap on seed 42 yields a 95\% interval $[0.006, 0.138]$---excluding 0, but wide, as Limited occupancy requires. DKT $\Delta\mathrm{FM}$ is positive on only three of five seeds and is not treated as a five-run finding. On XES3G5M, SimpleKT $\Delta\mathrm{FM}$ is negative on all five seeds, but $\Delta\mathrm{Miss}$ is positive on all five (mean $+0.112$): among actual incorrect answers, the system still advances more often on sparse than on dense KCs. A flat ECE therefore does not imply a flat miss rate.

\begin{table}[!t]
\caption{Gate robustness at $\tau=0.7$ on ASSISTments 2012 (five training seeds; four unique partitions). GKT/CL4KT remain seed 42 only.}
\label{tab:five}
\centering
\begin{tabular}{lcccc}
\toprule
Model & Mean $\Delta$FM & SD & Seeds $>0$ & Mean sparse $N$ \\
\midrule
SimpleKT & 0.047 & 0.033 & 5/5 & 413 \\
DKT & 0.033 & 0.048 & 3/5 & 413 \\
\bottomrule
\end{tabular}
\end{table}

Coincidence of digits: SimpleKT dense $E[\mathrm{FM}]$ at $\tau=0.7$ on seed 42 is $0.113$, close to the four-partition dense ECE $0.114$. They are different quantities.

\section{Discussion}

\subsection{What a designer should take away}

If a learning platform uses KT probabilities only to sort students, Table~\ref{tab:overall} may suffice. If it uses $p$ to skip practice or declare mastery, Table~\ref{tab:overall} is the wrong dashboard. The ASSISTments SimpleKT cell is the cautionary case: competitive AUC, Limited-but-estimable sparse occupancy (about 19\% of KCs fall below the sparse threshold on fold 0), a monotonic ECE rise, and a gate that advances more dirty sparse attempts. The XES3G5M cell is the opposite caution: plenty of sparse mass and Reliable occupancy, yet SimpleKT ECE does not degrade---while miss rates still can. Junyi shows that the protocol can refuse a sparse claim when the bucket is empty.

Three observable conditions track when sparse-calibration claims are informative rather than obligatory (Table~\ref{tab:cond}): (1)~a non-empty low-frequency tail under the split actually used; (2)~enough sparse test events for at least a Limited-flag ECE; (3)~frequency--difficulty coupling that is not inverted. ASSISTments meets all three. That is a diagnostic pre-condition, not a claim that training frequency causes miscalibration.

\begin{table}[!t]
\caption{When a sparse-calibration claim is estimable (FINDINGS, not causes). Sparse mass: share of KCs with $f_{\mathrm{train}}<100$.}
\label{tab:cond}
\centering
\footnotesize
\setlength{\tabcolsep}{2.5pt}
\begin{tabular}{p{1.55cm}p{1.7cm}p{1.55cm}p{1.7cm}}
\toprule
Condition & ASSISTments & Junyi & XES3G5M \\
\midrule
Sparse mass & 18.9\% of KCs & 0\% & 22.5\% of KCs \\
Sparse $N$ & 415 (L) & empty & 2{,}010 (R) \\
Difficulty coupling & $\rho=-0.227$ (sparse harder) & $\rho=-0.416$ & $\rho=+0.087$ (weak, inverted) \\
SimpleKT ECE & $0.114\to 0.228$ & dense$\to$medium only & flat $0.114\to 0.125$ \\
\bottomrule
\end{tabular}
\end{table}

\subsection{What this paper does not show}

The gate is not a classroom policy and not an A/B test of remediation. GKT and the CL4KT adapter are single-fold, ASSISTments-only instantiations. Temporal results are a single corrected cutoff (seed 42), not a multi-cutoff variance estimate. IRT is not a fair ranking baseline on unseen learners. We did not tune $\tau$ on sparse events, and we do not recommend doing so in deployment without a separate validation design.

\subsection{Implications for information and education technology}

IJIET's audience includes operators of educational platforms, not only KT researchers. For that audience, the operational recommendation is narrow: (i)~log train-only KC frequency with every prediction; (ii)~report ECE and occupancy on dense vs.\ sparse slices before enabling a global mastery threshold; (iii)~if a threshold must be global, simulate FM and Miss on the sparse slice at the \emph{same} $\tau$ used in the product; (iv)~do not treat a population AUC win, or a graph/contrastive architecture, as automatically safer on the tail.

\section{Conclusion}

KT models that look adequate on AUC can still be poorly calibrated on rarely practiced skills, and a global mastery threshold can then produce a higher false-mastery rate on those skills. That pattern appears for SimpleKT on ASSISTments 2012 (ECE $0.114\to 0.228$, Limited $N=415$; seed-42 FM $0.196\to 0.268$; $\Delta\mathrm{FM}>0$ on 5/5 seeds). It does not appear as a universal law: Junyi has no learner-based sparse stratum, and XES3G5M SimpleKT ECE is flat. Sparse-concept diagnostics are therefore conditionally important. The simulation is a decision-error check for educational-technology systems, not a claim of a new KT model and not a classroom intervention.

\section*{Conflict of Interest}

The authors declare no conflict of interest.

\section*{Author Contributions}

K.-T. Nguyen led the diagnostic design, experiments, and manuscript. T. Dao Minh and D. Nguyen Tien contributed data processing and baseline runs. C. T. Nguyen contributed methodological review. V.-H. Nguyen supervised the study and revised the manuscript. All authors approved the final version.

\section*{Acknowledgment}

Anonymized code and prediction exports are available for peer review; a public repository will be released upon acceptance.

\bibliographystyle{IEEEtran}
\bibliography{../paper/references}

\end{document}
